\section{Aplicación empresarial}

El análisis financiero permite convertir resultados numéricos en decisiones concretas. Si las ventas crecen, pero el costo de ventas aumenta a mayor velocidad, la empresa debe revisar precios, proveedores, procesos productivos y eficiencia operativa. Si el patrimonio pierde participación relativa frente al total activo, también conviene evaluar la estructura de financiamiento.

En la gestión empresarial, el análisis horizontal ayuda a detectar tendencias, mientras que el análisis vertical permite observar la composición económica y financiera. Ambos métodos pueden usarse en presupuestos, control de costos, evaluación de rentabilidad, decisiones de inversión y seguimiento de metas.

\section{Aplicación del análisis financiero en Ingeniería de Sistemas}

La Contabilidad Empresarial y la Ingeniería de Sistemas se relacionan mediante los sistemas de información que registran, procesan y reportan datos económicos. En una empresa, las operaciones de ventas, compras, bancos, inventarios y contabilidad pueden integrarse en un ERP, permitiendo que la información contable sea consultada y transformada en reportes financieros.

Las bases de datos permiten almacenar ventas, costos, activos, pasivos, patrimonio, compras e inventarios de manera estructurada. A partir de estas tablas, un analista o ingeniero de sistemas puede construir consultas que agrupen información por periodo, cuenta contable, área, proyecto o centro de costo.

El análisis horizontal puede automatizarse mediante hojas de cálculo, SQL, Python, Power BI o módulos de ERP que comparen el periodo actual con el periodo anterior. Del mismo modo, el análisis vertical puede calcularse dividiendo automáticamente cada partida entre una base común, como ventas netas o total de activos.

Los dashboards financieros permiten presentar resultados de forma visual, mostrando variaciones, márgenes, estructura de costos, liquidez, deuda y rentabilidad. Estos tableros pueden incluir alertas cuando los costos crecen más rápido que las ventas, cuando el margen bruto disminuye o cuando el endeudamiento supera un límite definido.

Como ejemplo laboral, una empresa tecnológica puede monitorear ventas por licencias, costos de infraestructura cloud, gastos administrativos, soporte técnico y utilidad neta. Si el gasto cloud pasa de S/ 8,000.00 a S/ 12,000.00, el análisis horizontal mide el incremento; si ese gasto representa un porcentaje creciente de las ventas, el análisis vertical ayuda a evaluar si la estructura de costos sigue siendo sostenible.

Por ello, el análisis financiero no solo es una herramienta contable. También es un campo de aplicación para la Ingeniería de Sistemas, porque permite diseñar reportes, automatizar indicadores y apoyar decisiones gerenciales con información oportuna, trazable y útil.

\subsection{Aplicación tecnológica del análisis financiero}

Una solución tecnológica de análisis financiero puede organizarse como un flujo: ERP, base de datos, proceso ETL, cálculo de indicadores, dashboard y alerta gerencial. En este flujo, el ERP registra operaciones; la base de datos almacena transacciones; el ETL limpia y organiza la información; el cálculo financiero genera variaciones, porcentajes y ratios; y el dashboard presenta resultados visuales para la toma de decisiones.

Herramientas como Excel, SQL, Python, Power BI o sistemas web permiten automatizar reportes y reducir errores de cálculo. Esta automatización facilita que la gerencia revise tendencias, márgenes, deuda, liquidez y rentabilidad con información actualizada.

\section{Conclusiones}

\begin{itemize}
  \item El análisis financiero convierte estados financieros en información para decidir.
  \item El análisis horizontal permite observar tendencias y variaciones entre periodos.
  \item El análisis vertical permite estudiar la estructura porcentual de un estado financiero.
  \item En la Empresa ABC, las ventas crecen, pero el costo de ventas crece más rápido, reduciendo el margen bruto.
  \item Los resultados deben complementarse con ratios de liquidez, gestión, solvencia y rentabilidad.
  \item Para Ingeniería de Sistemas, estos análisis pueden automatizarse mediante ERP, bases de datos y dashboards.
\end{itemize}
